% Archived custom LaTeX/TikZ sources removed from Chapter 3 on 2026-06-26.
% They are not included by main.tex. Keep this file as recoverable source material.
% Source revision before removal: aa56e5c (after framing cleanup, before Chapter 3 rewrite).
% Current removal commit: 5be32f5.

%==============================================================================
% Embedding inversion TikZ figure
% Original label: fig:method-embedding-inversion-example
%==============================================================================
\begin{figure}[!htbp]
    \centering
    \resizebox{0.98\textwidth}{!}{%
    \begin{tikzpicture}[
        font=\scriptsize,
        >=Latex,
        lane/.style={draw=#1!60!black, fill=#1!5, very thick, rounded corners=6pt},
        card/.style={draw=#1!65!black, fill=white, very thick, rounded corners=4pt,
            align=center, text width=25mm, minimum height=16mm, inner sep=3.2pt},
        badge/.style={circle, draw=white, line width=0.8pt, fill=#1!18,
            minimum size=8.2mm, inner sep=0pt},
        lane tag/.style={draw=#1!65!black, fill=#1!12, very thick,
            rounded corners=4pt, align=center, text width=20mm,
            minimum height=12mm, inner sep=3pt, font=\scriptsize\bfseries},
        line/.style={-{Latex[length=2.2mm]}, very thick, draw=#1!65!black},
        turn/.style={-{Latex[length=2.2mm]}, very thick, draw=#1!65!black,
            rounded corners=7pt},
        optional/.style={draw=#1!65!black, densely dashed},
        pics/document/.style={code={
            \draw[fill=white, draw=black!65, rounded corners=0.3pt] (-0.15,-0.15) rectangle (0.14,0.18);
            \draw[fill=white, draw=black!65, rounded corners=0.3pt] (-0.07,-0.08) rectangle (0.22,0.25);
            \draw[draw=black!55, line width=0.35pt] (-0.02,0.15) -- (0.16,0.15);
            \draw[draw=black!55, line width=0.35pt] (-0.02,0.07) -- (0.16,0.07);
        }},
        pics/scale/.style={code={
            \draw[draw=black!65, line width=0.55pt] (0,-0.22) -- (0,0.18);
            \draw[draw=black!65, line width=0.55pt] (-0.20,0.08) -- (0.20,0.08);
            \draw[draw=black!65, line width=0.45pt] (-0.15,0.08) -- (-0.23,-0.07) -- (-0.07,-0.07) -- cycle;
            \draw[draw=black!65, line width=0.45pt] (0.15,0.08) -- (0.07,-0.07) -- (0.23,-0.07) -- cycle;
            \draw[draw=black!65, line width=0.55pt] (-0.10,-0.22) -- (0.10,-0.22);
        }},
        pics/feedback/.style={code={
            \draw[fill=white, draw=black!65, rounded corners=1pt] (-0.22,-0.08) rectangle (0.10,0.15);
            \draw[fill=white, draw=black!65, rounded corners=1pt] (-0.06,-0.17) rectangle (0.24,0.06);
            \draw[draw=black!55, line width=0.35pt] (-0.16,0.06) -- (0.04,0.06);
            \draw[draw=black!55, line width=0.35pt] (0.00,-0.03) -- (0.18,-0.03);
            \fill[darkyellow!80!black] (0.19,0.14) circle (0.035);
        }},
        pics/prompt/.style={code={
            \draw[fill=white, draw=black!65, rounded corners=1pt] (-0.20,-0.17) rectangle (0.18,0.17);
            \draw[draw=black!55, line width=0.35pt] (-0.13,0.08) -- (0.11,0.08);
            \draw[draw=black!55, line width=0.35pt] (-0.13,0.00) -- (0.08,0.00);
            \draw[draw=orange!80!black, line width=0.75pt] (-0.06,-0.10) -- (0.16,0.12);
            \filldraw[fill=orange!45, draw=orange!80!black] (0.16,0.12) circle (0.025);
        }},
        pics/lens/.style={code={
            \draw[draw=black!65, line width=0.55pt] (-0.03,0.04) circle (0.13);
            \draw[draw=black!65, line width=0.55pt] (0.07,-0.06) -- (0.22,-0.21);
            \draw[draw=darkgreen!70!black, line width=0.7pt] (-0.09,0.03) -- (-0.02,-0.05) -- (0.10,0.10);
        }},
        pics/archive/.style={code={
            \draw[fill=white, draw=black!65, rounded corners=1pt] (-0.22,-0.15) rectangle (0.22,0.16);
            \draw[draw=black!55, line width=0.35pt] (-0.14,0.07) -- (0.14,0.07);
            \draw[draw=black!55, line width=0.35pt] (-0.14,-0.01) -- (0.10,-0.01);
            \draw[draw=darkgreen!70!black, line width=0.65pt] (-0.10,-0.09) -- (-0.03,-0.16) -- (0.14,0.01);
        }}
    ]
        \node[lane=orange, minimum width=16.4cm, minimum height=2.0cm, anchor=west] at (0,0) {};
        \node[lane=green, minimum width=16.4cm, minimum height=2.0cm, anchor=west] at (0,-2.35) {};

        \node[lane tag=orange, text=orange!70!black] at (1.15,0) {Validation\\feedback};
        \node[lane tag=green, text=green!45!black] at (1.15,-2.35) {Prompt\\selection};

        \node[card=orange] (example) at (3.50,0) {\textbf{Evaluation example}\\dialogue, response\\human score};
        \node[card=orange] (judge) at (6.40,0) {\textbf{Base judge}\\Qwen2.5-7B\\predicted score};
        \node[card=orange] (reflection) at (9.30,0) {\textbf{Reflection data}\\metric error\\optional PPL/NLA};
        \node[card=orange, optional=orange] (aux) at (12.20,0) {\textbf{Aux judge}\\optional\\feedback compression};
        \node[card=orange] (proposer) at (15.10,0) {\textbf{Proposer}\\Qwen35B\\candidate edit};

        \node[card=green] (candidate) at (15.10,-2.35) {\textbf{Candidate prompt}\\new judge\\instructions};
        \node[card=green] (validation) at (11.60,-2.35) {\textbf{Validation}\\score candidate\\on held-out val};
        \node[card=green] (archive) at (8.10,-2.35) {\textbf{Prompt pool}\\keep if useful\\or discard};
        \node[card=green] (final) at (4.60,-2.35) {\textbf{Final test}\\run once after\\selection};

        \node[badge=orange] at ([yshift=2.1mm]example.north) {};
        \pic[scale=0.78] at ([yshift=2.1mm]example.north) {document};
        \node[badge=orange] at ([yshift=2.1mm]judge.north) {};
        \pic[scale=0.78] at ([yshift=2.1mm]judge.north) {scale};
        \node[badge=orange] at ([yshift=2.1mm]reflection.north) {};
        \pic[scale=0.78] at ([yshift=2.1mm]reflection.north) {feedback};
        \node[badge=orange] at ([yshift=2.1mm]aux.north) {};
        \pic[scale=0.78] at ([yshift=2.1mm]aux.north) {lens};
        \node[badge=orange] at ([yshift=2.1mm]proposer.north) {};
        \pic[scale=0.78] at ([yshift=2.1mm]proposer.north) {prompt};

        \node[badge=green] at ([yshift=2.1mm]candidate.north) {};
        \pic[scale=0.78] at ([yshift=2.1mm]candidate.north) {prompt};
        \node[badge=green] at ([yshift=2.1mm]validation.north) {};
        \pic[scale=0.78] at ([yshift=2.1mm]validation.north) {scale};
        \node[badge=green] at ([yshift=2.1mm]archive.north) {};
        \pic[scale=0.78] at ([yshift=2.1mm]archive.north) {archive};
        \node[badge=green] at ([yshift=2.1mm]final.north) {};
        \pic[scale=0.78] at ([yshift=2.1mm]final.north) {lens};

        \draw[line=orange] (example.east) -- (judge.west);
        \draw[line=orange] (judge.east) -- (reflection.west);
        \draw[line=orange] (reflection.east) -- (aux.west);
        \draw[line=orange] (aux.east) -- (proposer.west);
        \draw[turn=orange] (proposer.south) -- ++(0,-0.18) -| ([xshift=7mm]candidate.north);
        \draw[line=green] (candidate.west) -- (validation.east);
        \draw[line=green] (validation.west) -- (archive.east);
        \draw[line=green] (archive.west) -- (final.east);
    \end{tikzpicture}}
    \caption{Detailed GEPA feedback flow used in the prompt-optimization track.
    Reflection data are produced on validation examples, optionally enriched with
    perplexity, NLA verbalizations, or an auxiliary judge, and then given to the
    proposer to generate candidate prompts. Final metrics are computed only after
    prompt selection; test examples do not drive proposal or validation-time
    selection.}
    \label{fig:method-gepa-feedback-flow}
\end{figure}

\FloatBarrier

\section{Canonical Semantic-Fidelity Dataset}
\label{sec:method-canonical-dataset}

The thesis uses a canonical semantic-fidelity dataset to stress latent-to-text
methods. The dataset is not intended to be a large benchmark by itself. Its role
is to provide controlled examples where a reconstruction can look plausible
while changing the logical meaning of the input.

Each row is represented with stable metadata. At minimum, a row contains a
stable identifier, a source block, a phenomenon label, an input sentence, an
optional paired sentence, and split metadata:

\begin{itemize}
    \item the stable identifier makes outputs from different branches joinable;
    \item the source block records whether the row is a standard control,
    negation/polarity item, or commonsense/counterfactual item;
    \item the input sentence is the text given to latent-to-text methods;
    \item the paired sentence is used when the phenomenon is defined by a
    contrast, such as a negated sentence and its non-negated counterpart;
    \item the split metadata prevents train/validation/test leakage in later
    experiments.
\end{itemize}

\begin{table}[!htb]
    \centering
    \small
    \begin{adjustbox}{width=\textwidth}
    \begin{tabular}{p{0.10\linewidth} p{0.18\linewidth} p{0.29\linewidth} p{0.29\linewidth} p{0.18\linewidth}}
        \toprule
        \textbf{Block} & \textbf{Phenomenon} & \textbf{Input text} & \textbf{Paired text} & \textbf{Expected issue} \\
        \midrule
        A & Standard control &
        A simple factual or descriptive sentence. &
        Not required. &
        Checks whether a method works on ordinary inputs. \\
        B & Negation and polarity &
        The treatment did not improve the patient's condition. &
        The treatment improved the patient's condition. &
        Small lexical changes reverse the truth condition. \\
        C & Commonsense or counterfactual content &
        A fragile glass fell onto concrete and did not break. &
        A fragile glass fell onto concrete and broke. &
        Plausible reconstructions may normalize the sentence toward commonsense. \\
        \bottomrule
    \end{tabular}
    \end{adjustbox}
    \caption{Canonical semantic-fidelity dataset structure. The examples are
    illustrative; the final dataset rows are specified in the experimental
    setup.}
    \label{tab:canonical-dataset-blocks}
\end{table}

\TableRef{tab:canonical-dataset-blocks} shows the logic of the dataset. Block A
acts as a control: a method should not be judged only on adversarial or unusual
sentences. Block B focuses on negation and polarity, where a single word can
change the truth value of the statement. Block C focuses on examples that are
counterfactual or contrary to common expectations. These examples are useful
because language models tend to prefer plausible continuations. A latent-to-text
method that silently normalizes an unusual sentence can therefore produce text
that is fluent and semantically close in a loose sense, but wrong with respect
to the property under investigation.

The dataset is used as a common reference object across branches. In the SIPIT
branch, rows can be tokenized into fixed prompt windows. In the NLA branch, the
same rows define which activation positions should be verbalized. In the
embedding-inversion branch, they provide diagnostic inputs for observing whether
reconstructions preserve logical content.

\section{Embedding-Inversion Diagnostic Branch}
\label{sec:method-embedding-inversion}

\subsection{Task Setting}
\label{subsec:method-embedding-task}

The first latent-to-text branch concerns embedding inversion. In this setting,
the input to the inversion method is a fixed-size text embedding, and the output
is a textual reconstruction. The thesis uses this branch as a diagnostic track:
it investigates how embedding-inversion systems behave on semantically fragile
inputs, but it does not rely on this branch alone for the final claim.
\FigureRef{fig:method-embedding-inversion-example} shows the practical
operation of the script-level pipeline on one local qualitative artifact: the
system first encodes the sentence, then decodes from the embedding alone, and
finally exposes whether the recovered text still expresses the intended
predicate and negation.

\begin{figure}[!htbp]
    \centering
    \resizebox{0.98\linewidth}{!}{%
    \begin{tikzpicture}[
        font=\scriptsize,
        >=Latex,
        stage/.style={draw=#1!65!black, fill=white, very thick, rounded corners=4pt,
            align=center, text width=22mm, minimum height=22mm, inner sep=3.0pt},
        wide stage/.style={draw=#1!65!black, fill=white, very thick, rounded corners=4pt,
            align=left, text width=95mm, minimum height=18mm, inner sep=4pt},
        note/.style={draw=#1!65!black, fill=#1!8, very thick, rounded corners=4pt,
            align=left, text width=105mm, minimum height=16mm, inner sep=4pt},
        lane/.style={draw=#1!55!black, fill=#1!5, very thick, rounded corners=6pt},
        badge/.style={circle, draw=white, line width=0.8pt, fill=#1!18,
            minimum size=8.5mm, inner sep=0pt},
        line/.style={-{Latex[length=2.2mm]}, very thick, draw=#1!65!black,
            shorten >=0.8mm, shorten <=0.8mm},
        small label/.style={font=\scriptsize\bfseries, text=#1!70!black},
        pics/doc/.style={code={
            \draw[fill=white, draw=black!65, rounded corners=0.3pt] (-0.17,-0.17) rectangle (0.18,0.20);
            \draw[draw=black!55, line width=0.35pt] (-0.10,0.10) -- (0.11,0.10);
            \draw[draw=black!55, line width=0.35pt] (-0.10,0.02) -- (0.14,0.02);
            \draw[draw=black!55, line width=0.35pt] (-0.10,-0.06) -- (0.08,-0.06);
        }},
        pics/vector/.style={code={
            \foreach \x/\h/\c in {-0.21/0.18/blue,-0.11/0.09/green,0.00/0.23/orange,0.11/0.13/lightviolet,0.22/0.05/darkyellow}
                \draw[fill=\c!35, draw=\c!70!black, rounded corners=0.5pt] (\x,-0.18) rectangle ++(0.065,\h);
            \draw[draw=black!45, line width=0.35pt] (-0.25,-0.18) -- (0.30,-0.18);
        }},
        pics/mask/.style={code={
            \foreach \x in {-0.24,-0.12,0,0.12,0.24}
                \draw[fill=black!10, draw=black!55, rounded corners=1pt] (\x-0.045,-0.06) rectangle (\x+0.045,0.06);
            \node[font=\tiny] at (0,0) {?};
        }},
        pics/decoder/.style={code={
            \draw[draw=black!60, line width=0.55pt] (-0.22,-0.14) rectangle (0.22,0.14);
            \draw[draw=black!50, line width=0.4pt] (-0.14,-0.21) -- (-0.05,-0.14);
            \draw[draw=black!50, line width=0.4pt] (0.14,-0.21) -- (0.05,-0.14);
            \draw[draw=blue!70!black, line width=0.6pt] (-0.15,0.03) -- (-0.02,0.10) -- (0.14,-0.04);
            \fill[orange!70!black] (0.15,-0.04) circle (0.03);
        }},
        pics/warn/.style={code={
            \draw[fill=orange!20, draw=orange!75!black, line width=0.55pt] (0,0.22) -- (-0.22,-0.18) -- (0.22,-0.18) -- cycle;
            \draw[draw=orange!80!black, line width=0.55pt] (0,0.08) -- (0,-0.06);
            \fill[orange!80!black] (0,-0.12) circle (0.018);
        }}
    ]
        \node[lane=blue, minimum width=11.7cm, minimum height=5.65cm, anchor=west] at (0,-1.62) {};

        \node[stage=blue] (source) at (1.35,0.00) {\textbf{Input text}\\
            \quotes{A man not playing guitar on stage.}};
        \node[stage=green] (embed) at (4.25,0.00) {\textbf{Encoder}\\
            Jina/Qwen3 demo\\embedding vector};
        \node[stage=orange] (masked) at (7.15,0.00) {\textbf{Initial decoder state}\\
            all text positions\\masked};
        \node[stage=lightviolet] (decode) at (10.05,0.00) {\textbf{Conditional MDLM}\\
            iterative\\unmasking};
        \node[wide stage=darkgreen] (output) at (5.85,-1.95) {\textbf{Recovered text}\\
            \textit{A man on a stage guitar playing a guitar on a flat bench.}\\
            \textit{He is not a stage. A man is playing the role of a music player.}};

        \node[badge=blue] at ([yshift=2.1mm]source.north) {};
        \pic[scale=0.78] at ([yshift=2.1mm]source.north) {doc};
        \node[badge=green] at ([yshift=2.1mm]embed.north) {};
        \pic[scale=0.78] at ([yshift=2.1mm]embed.north) {vector};
        \node[badge=orange] at ([yshift=2.1mm]masked.north) {};
        \pic[scale=0.78] at ([yshift=2.1mm]masked.north) {mask};
        \node[badge=lightviolet] at ([yshift=2.1mm]decode.north) {};
        \pic[scale=0.78] at ([yshift=2.1mm]decode.north) {decoder};
        \draw[line=blue] (source.east) -- (embed.west);
        \draw[line=green] (embed.east) -- (masked.west);
        \draw[line=orange] (masked.east) -- (decode.west);
        \draw[line=lightviolet] (decode.south) -- ([xshift=39mm]output.north);

        \node[note=orange] (check) at (5.85,-3.68) {\textbf{Semantic check}\\
            The word \quotes{not} appears, but it attaches to a different
            predicate (\quotes{not a stage}); the main event drifts back toward
            \quotes{playing guitar}.};
        \node[badge=orange] at ([xshift=-6.0mm]check.west) {};
        \pic[scale=0.68] at ([xshift=-6.0mm]check.west) {warn};
        \draw[line=orange] (output.south) -- (check.north);
    \end{tikzpicture}}
    \caption[Embedding-inversion script example.]{Concrete embedding-inversion example from the local
    \texttt{eval\_metrics.py} artifact. The example was generated through the
    public Jina demo encode/decode path and is used here only to show what the
    script does in practice: from a sentence it obtains an embedding, decodes
    from that embedding, and then exposes a semantic-fidelity issue.}
    \label{fig:method-embedding-inversion-example}
\end{figure}

%==============================================================================
% SIPIT terminal-style listing
% Original label: lst:sipit-terminal-example
%==============================================================================


%==============================================================================
% Standalone NLA TikZ figure
% Original label: fig:method-nla-standalone
%==============================================================================
\begin{figure}[!htb]
    \centering
    \resizebox{0.98\linewidth}{!}{%
    \begin{tikzpicture}[
        font=\scriptsize,
        >=Latex,
        lane/.style={draw=#1!60!black, fill=#1!5, very thick, rounded corners=6pt},
        card/.style={draw=#1!65!black, fill=white, very thick, rounded corners=4pt,
            align=center, text width=25mm, minimum height=16mm, inner sep=3.2pt},
        badge/.style={circle, draw=white, line width=0.8pt, fill=#1!18,
            minimum size=8.2mm, inner sep=0pt},
        lane title/.style={anchor=west, font=\small\bfseries, align=left},
        line/.style={-{Latex[length=2.2mm]}, very thick, draw=#1!65!black},
        pics/document/.style={code={
            \draw[fill=white, draw=black!65, rounded corners=0.3pt] (-0.15,-0.15) rectangle (0.14,0.18);
            \draw[fill=white, draw=black!65, rounded corners=0.3pt] (-0.07,-0.08) rectangle (0.22,0.25);
            \draw[draw=black!55, line width=0.35pt] (-0.02,0.15) -- (0.16,0.15);
            \draw[draw=black!55, line width=0.35pt] (-0.02,0.07) -- (0.16,0.07);
        }},
        pics/network/.style={code={
            \draw[draw=black!55, line width=0.45pt] (-0.18,0.13) -- (0.02,0.00) -- (0.20,0.13);
            \draw[draw=black!55, line width=0.45pt] (-0.18,-0.13) -- (0.02,0.00) -- (0.20,-0.13);
            \foreach \x/\y in {-0.18/0.13,-0.18/-0.13,0.02/0,0.20/0.13,0.20/-0.13}
                \filldraw[fill=white, draw=black!65] (\x,\y) circle (0.045);
        }},
        pics/tokens/.style={code={
            \foreach \x/\c in {-0.20/blue, -0.07/orange, 0.06/green, 0.19/lightviolet}
                \filldraw[fill=\c!25, draw=\c!70!black] (\x,0) circle (0.055);
            \draw[draw=black!45, line width=0.4pt] (-0.145,0) -- (-0.125,0);
            \draw[draw=black!45, line width=0.4pt] (-0.015,0) -- (0.005,0);
            \draw[draw=black!45, line width=0.4pt] (0.115,0) -- (0.135,0);
        }},
        pics/lens/.style={code={
            \draw[draw=black!65, line width=0.55pt] (-0.03,0.04) circle (0.13);
            \draw[draw=black!65, line width=0.55pt] (0.07,-0.06) -- (0.22,-0.21);
            \draw[draw=darkgreen!70!black, line width=0.7pt] (-0.09,0.03) -- (-0.02,-0.05) -- (0.10,0.10);
        }},
        pics/archive/.style={code={
            \draw[fill=white, draw=black!65, rounded corners=1pt] (-0.22,-0.15) rectangle (0.22,0.16);
            \draw[draw=black!55, line width=0.35pt] (-0.14,0.07) -- (0.14,0.07);
            \draw[draw=black!55, line width=0.35pt] (-0.14,-0.01) -- (0.10,-0.01);
            \draw[draw=darkgreen!70!black, line width=0.65pt] (-0.10,-0.09) -- (-0.03,-0.16) -- (0.14,0.01);
        }}
    ]
        \node[lane=blue, minimum width=9.25cm, minimum height=2.65cm, anchor=west] at (0,0) {};
        \node[lane=green, minimum width=6.50cm, minimum height=2.65cm, anchor=west] at (9.60,0) {};
        \node[lane title, text=blue!55!black] at (0.35,1.48) {Activation extraction};
        \node[lane title, text=green!45!black] at (9.95,1.48) {Natural-language readout};

        \node[card=blue] (input) at (1.90,-0.18) {\textbf{Input example}\\source text\\selected token};
        \node[card=blue] (base) at (5.05,-0.18) {\textbf{Base LM}\\Qwen2.5-7B\\forward pass};
        \node[card=blue] (activation) at (8.20,-0.18) {\textbf{Layer-20 vector}\\residual stream\\token position};
        \node[card=green] (av) at (11.25,-0.18) {\textbf{NLA AV}\\checkpoint\\vector to text};
        \node[card=green] (record) at (14.30,-0.18) {\textbf{Saved record}\\verbalization\\token/layer metadata};

        \node[badge=blue] at ([yshift=2.1mm]input.north) {};
        \pic[scale=0.78] at ([yshift=2.1mm]input.north) {document};
        \node[badge=blue] at ([yshift=2.1mm]base.north) {};
        \pic[scale=0.78] at ([yshift=2.1mm]base.north) {network};
        \node[badge=blue] at ([yshift=2.1mm]activation.north) {};
        \pic[scale=0.78] at ([yshift=2.1mm]activation.north) {tokens};
        \node[badge=green] at ([yshift=2.1mm]av.north) {};
        \pic[scale=0.78] at ([yshift=2.1mm]av.north) {lens};
        \node[badge=green] at ([yshift=2.1mm]record.north) {};
        \pic[scale=0.78] at ([yshift=2.1mm]record.north) {archive};

        \draw[line=blue] (input.east) -- (base.west);
        \draw[line=blue] (base.east) -- (activation.west);
        \draw[line=green] (activation.east) -- (av.west);
        \draw[line=green] (av.east) -- (record.west);
    \end{tikzpicture}}
    \caption{Standalone NLA extraction pipeline used before integrating
    verbalizations into GEPA feedback. The left side extracts a residual-stream
    vector from the base model at a selected token position; the right side turns
    that vector into a natural-language record with token, layer, and source
    metadata.}
    \label{fig:method-nla-standalone}
\end{figure}

%==============================================================================
% NLA verbalization example table
% Original label: tab:nla-standalone-example
%==============================================================================
\begin{table}[!htb]
    \centering
    \scriptsize
    \begin{adjustbox}{width=\textwidth}
    \begin{tabular}{p{0.22\linewidth} p{0.16\linewidth} p{0.15\linewidth} p{0.39\linewidth}}
        \toprule
        \textbf{Artifact field} & \textbf{Position} & \textbf{Token} & \textbf{Verbalization} \\
        \midrule
        \shortstack[l]{\texttt{context\_016}\\\texttt{response\_01}} &
        candidate middle &
        \texttt{warriors} &
        \quotes{team} or \quotes{and basketball} or \quotes{but my favorite sports team} \\
        \shortstack[l]{\texttt{context\_020}\\\texttt{response\_00}} &
        candidate middle &
        \texttt{made} &
        or \quotes{a tiny amount of money,} completing the absurdity about box office failure \\
        \shortstack[l]{\texttt{context\_011}\\\texttt{response\_03}} &
        candidate middle &
        \texttt{human} &
        the parallel phrase about human diseases, continuing the specific disease-correlation example \\
        \shortstack[l]{\texttt{context\_054}\\\texttt{response\_05}} &
        candidate middle &
        \texttt{need} &
        or \quotes{a friend to fill my streaming void,} completing the informal request \\
        \bottomrule
    \end{tabular}
    \end{adjustbox}
    \caption{Example NLA verbalizations from saved Topical-Chat artifacts. Each
    row reports one selected token from a candidate response and the corresponding
    Layer-20 activation readout. The selected examples show sports-team semantics,
    numerical box-office contradiction, factual biological relation, and
    conversational intent.}
    \label{tab:nla-standalone-example}
\end{table}
